\chapter{Introduction}
In recent years, quantum algorithms have garnered significant attention for their potential to disrupt traditional cryptographic systems. The security of practically all important communication channels, including the internet as we know it, is ensured by classical cryptographic protocols based on the hardness of certain mathematical problems. With the rise of quantum technology, new algorithms have been found that exploit quantum mechanical principles, such as superposition and entanglement, to solve these problems exponentially faster than classical algorithms. Still, the implementation of quantum algorithms for cryptographic purposes continues to be largely theoretical. A quantum computer capable of breaking the most-used algorithms today is yet to be constructed, as there are many practical limitations to doing so. Thus, practical quantum computers, with thousands of qubits and low error rates, remain an area of active research; there is considerable anticipation surrounding the eventual realization of the theoretical promises associated with quantum computing. As we move into the post-quantum era, it will become increasingly important to assess our cryptographic protocols in light of the most recent developments in the field in order to protect critical infrastructure and assets.

Brute-force search is a technique that may be employed to break cryptographic systems with \emph{symmetric} keys, which means the same keys are used for encryption and decryption \cite{symmetric}, as opposed to \emph{asymmetric} algorithms, which use separate encryption and decryption keys. Performing such a search involves quite simply checking every single possible key option until the right key is found, an inordinately time-consuming process if the key space is large enough. Our current protocols ensure security against these searches by having key lengths large enough that it is considered computationally infeasible for a classical computer to identify one in any reasonable amount of time. A closer look at the popular symmetric-key algorithm \emph{Advanced Encryption Standard} (AES) \cite{aes}, which encrypts data by applying a secret key through a series of mathematical transformations, exemplifies this: It exists in versions that use key sizes of 128, 192 and 256 bits, resulting in search spaces of $2^{128}$, $2^{192}$ and $2^{256}$ possible keys respectively.

A relevant quantum algorithm here is \emph{Grover's algorithm} \cite{Nielsen_Chuang_2010}, in which a procedure is given for a quantum computer to perform a search much faster than brute-force search can be done on classical computers, thereby potentially threatening the security of AES. As we will see, Grover's algorithm offers a quadratic improvement over classical algorithms, meaning the number of operations needed when we use this quantum algorithm is on the order of the square-root of the number of operations required when solving the same problem classically. However, finding the exact runtime is not as straightforward as the asymptotic estimates, because there are constant multiplicative factors hidden within the asymptotic notation that in practice determine the true resource cost of each operation. Asymptotic estimates also do not take into account what kind of parallelization may be possible. Analyzing the true cost of using Grover's algorithm is therefore a nontrivial problem that warrants further exploration, which constitutes the core of this project.

Building a quantum machine in practice is no easy feat. The root of the issue is that quantum information is extremely fragile. The qubits that store information are unstable and very vulnerable to noise that changes the value the qubit holds. We therefore wish to construct circuits of multiple qubits that work together to protect the information from errors, such that they form a single logical qubit that can have its errors located and corrected, making them more stable. This approach, though useful, will often require large numbers of qubits, which poses an engineering challenge to keep them all as stable as possible, facilitate connections between them, and perform computations on the logical qubits rather than on the physical qubits individually.

Jaques considers this in his 2024 invited talk about Grover's algorithm as a quantum attack on AES \cite{jaques2024ches_presentation}, in which he argues that due to the error correction models and gate sets currently in widespread use, a physical implementation of Grover's algorithm will never see the light of day on anything resembling current quantum architectures. His argument rests on three pillars: the enormous physical-to-logical qubit ratio imposed by current error-correcting codes, the prohibitive cost of fault-tolerant gate implementation under those codes, and the poor parallelism of Grover's algorithm itself. Since the first two pillars are explicitly tied to the surface code, the dominant error-correcting architecture in current fault-tolerant designs, a natural question arises: Does his conclusion survive under a more favorable architecture? This thesis investigates that question by examining what Jaques' analysis looks like under the assumption of the Bivariate Bicycle (BB) code family \cite{Bravyi2024}, a recent class of quantum low-density parity-check (LDPC) codes that offer substantially better qubit efficiency than the surface code.

A central challenge in this investigation is that quantum systems are notoriously difficult to simulate classically. If we could simulate quantum systems efficiently in general, it would beg the question why we need to realize them at all. This limits how directly we can study fault-tolerant implementations of Grover's algorithm: Even the smallest nontrivial instances quickly exceed the memory available on consumer machines. This thesis navigates that barrier using the stabilizer formalism \cite{paritycheckstabilizerwiki,pesah2023_stabilizer_trilogy1} as a computational tool, which enables efficient classical simulation of the error-correcting codes and logical gate operations relevant to our analysis, without requiring full statevector simulation.

The specific contributions of this thesis are as follows: We verify that fault-tolerant logical computation in the Steane code works as intended by implementing a complete instance of Grover's algorithm on two error-corrected logical qubits under depolarizing noise. We then compare the BB code against the surface code at the level of quantum memory and transversal gate noise using the \texttt{PanQEC} simulation framework, and interpret the results of \textcite{Xu_2026} on BB code magic state distillation factories in the context of Jaques' resource estimates. Together, these results allow us to assess which of Jaques' three pillars are weakened by adopting the BB code, and whether any are broken.

% % What am I going to do?\\
% As a starting point for this project, I aim to recreate the work of Samuel Jacues, as laid out in his 2024 talk. Jaques takes a very pessimistic view on the potential use cases for Grover's algorithm to break AES, by which I mean he believes these use cases are entirely nonexistent. His argument relies on a series of assumptions for this to be the case. In the beginning stages of this project, I will tackle these assumptions case by case and see if it changes the conclusion at all. Along the way, we will run into a number of topics that I am not really familiar with, and we will have to acquire some knowledge to say if the arguments presented by Jacques still hold in these alternative cases.


% \section{Quantum Computing Prerequisities and Notation}
% \supervisorsinline{Do I include this section?}
% \todo[inline]{Rewrite and formalize, if it is to be included at all.}
% We assume you know basic Dirac notation, maybe say that n is the number of qubits and N is $2^n$?

% Hadamard gate, superposition

% Pauli matrices are the matrices I, X, Y and Z, defined as...
% Sometimes $\sigma$ notation is used, but this text uses only the I, X, Y, Z notation (this is not true though, sigma notation shows up in surface code section about stabilizers).

% The probability of measuring a particular component of a state vector in superposition is its complex scalar squared: ``As we know, in a quantum state vector, the amplitude of any component squared represents the probability of seeing that component after measurement.''
\section{Thesis Structure}\label{sec:structure}

The remainder of this thesis is organized as follows: \Cref{chap:background} develops the theoretical background, covering Grover's algorithm, quantum error correction, the surface and BB codes, and logical gates in these codes. \Cref{chap:exp_and_analysis} describes the experiments: a fault-tolerant implementation of Grover's algorithm on the 7-qubit code as a proof-of-concept, followed by \texttt{PanQEC}-based comparisons of the BB code and surface code at the level of quantum memory and transversal gate noise, and magic state distillation analysis drawing on \textcite{Xu_2026}. \Cref{chap:conclusion} synthesizes the findings and returns to Jaques' three pillars to assess which survive under the BB code architecture. All code experiments and simulations can be found on GitHub\footnote{\url{https://github.com/emmastorberg/CS5960PHYS.git}}.


\section{Quantum Computing Prerequisites and Notation}\label{sec:prerequisites}
%\paragraph{Notation.} 
We write $\mathbb{I}_m$ for the $m \times m$ identity matrix. The set of complex $m \times n$ matrices is denoted $\mathbb{C}^{m \times n}$, and $\mathbb{R}$ and $\mathbb{C}$ denote the real and complex numbers respectively.

Throughout this text, $n$ denotes the number of qubits in a system, and we write $N = 2^n$ for the dimension of the corresponding Hilbert space. We assume familiarity with Dirac notation: A quantum state is written as a ket $\ket{\psi} \in \mathbb{C}^N$, its dual as a bra $\bra{\psi}$, and the inner product of two states as $\bra{\phi}\ket{\psi}$. Dirac notation is occasionally used outside of its quantum mechanical context as a 
convenient algebraic shorthand: $\ket{i}\bra{j}$ denotes the outer product of two basis vectors, defining an $N \times N$ matrix with a 1 in position $(i,j)$ and zeroes elsewhere, with zero-based indices throughout.

A single qubit occupies a two-dimensional Hilbert space spanned by the computational basis states $\ket{0}$ and $\ket{1}$. A general single-qubit state is a superposition
\[
    \ket{\psi} = \alpha\ket{0} + \beta\ket{1}, \qquad \alpha, \beta \in \mathbb{C},
    \quad |\alpha|^2 + |\beta|^2 = 1.
\]
More generally, for an $n$-qubit system, the $2^n$ computational basis states $\ket{x}$, $x \in \{0,1\}^n$, span the full Hilbert space, and the state of the system is a superposition $\ket{\psi} = \sum_{x} c_x \ket{x}$ with $c_x \in \mathbb{C}$. The probability of obtaining outcome $x$ upon measurement in the computational basis is $|c_x|^2$, so the normalization condition $\sum_x |c_x|^2 = 1$ ensures that the probabilities sum to 1.

Quantum gates are represented by unitary matrices acting on these state vectors. The four single-qubit gates used most frequently in this text are the Pauli operators:
\[
    \mathbb{I}_2 = \begin{pmatrix} 1 & 0 \\ 0 & 1 \end{pmatrix}, \quad
    X = \begin{pmatrix} 0 & 1 \\ 1 & 0 \end{pmatrix}, \quad
    Y = \begin{pmatrix} 0 & -i \\ i & 0 \end{pmatrix}, \quad
    Z = \begin{pmatrix} 1 & 0 \\ 0 & -1 \end{pmatrix}.
\]
These are sometimes written as $\sigma_0, \sigma_1, \sigma_2, \sigma_3$ (or $\sigma_x, \sigma_y, \sigma_z$ for the non-identity operators) in the physics literature. Where this text departs from the $\mathbb{I}_2, X, Y, Z$ convention, in particular in the discussion of stabilizer codes in \cref{sec:toriccode}, the notation is introduced locally.

Another gate of central importance is the single-qubit Hadamard gate:
\[
    H = \frac{1}{\sqrt{2}}\begin{pmatrix} 1 & 1 \\ 1 & -1 \end{pmatrix},
\]
which maps each computational basis state to an equal superposition: $H\ket{0} = \tfrac{1}{\sqrt{2}}(\ket{0}+\ket{1})$ and $H\ket{1} = \tfrac{1}{\sqrt{2}}(\ket{0}-\ket{1})$. Applying $H$ to every qubit of an $n$-qubit register initialized in $\ket{0}^{\otimes n}$ (the $n$-fold tensor product $\ket{0}^{\otimes n} = \ket{0}\otimes\cdots\otimes\ket{0}$) produces the uniform superposition $\frac{1}{\sqrt{N}}\sum_{x=0}^{N-1}\ket{x}$, a construction that appears repeatedly in quantum algorithms, including Grover's algorithm.